\documentclass[12pt]{article}
\usepackage[a4paper,margin=1in]{geometry}
\usepackage[utf8]{inputenc}
\usepackage[T1]{fontenc}
\usepackage{amsmath,amssymb}
\usepackage{textcomp}
\usepackage{parskip}
\begin{document}
\begin{center}
{\Large \textbf{2008 STEP II}}
\end{center}

\subsection*{Page 1}
\begin{flushleft}
\bigskip\textbf{Section A: Pure Mathematics}\par
\bigskip\textbf{Question 1.} A sequence of points (x , y ), (x , y ), . . . in the cartesian plane is generated by first choosing\par
1 1 2 2\par
\[(x , y ) then applying the rule, for n = 1, 2, . . .,\]
1 1\par
\[(x , y ) = (x^2 y^2 + a, 2x y + b + 2) ,\]
n+1 n+1 n - n n n\par
where a and b are given real constants.\par
\textbf{(i)} In the case a = 1 and b = 1, find the values of (x , y ) for which the sequence is\par
1 1\par
-\par
constant.\par
\textbf{(ii)} Given that (x , y ) = ( 1, 1), find the values of a and b for which the sequence has\par
1 1\par
-\par
period 2.\par
\bigskip\textbf{Question 2.} Let a be the coefficient of xn in the series expansion, in ascending powers of x, of\par
n\par
1 + x\par
,\par
\[(1 x)2(1 + x^2)\]
-\par
\[where x < 1 . Show, using partial fractions, that either a = n + 1 or a = n + 2 according\]
n n\par
| |\par
to the value of n.\par
11 000\par
Hence find a decimal approximation, to nine significant figures, for the fraction .\par
8181\par
[You are not required to justify the accuracy of your approximation.]\par
\bigskip\textbf{Question 3.} (i) Find the coordinates of the turning points of the curve y = 27x3 27x2 + 4. Sketch the\par
-\par
\[curve and deduce that x^2(1 x)  4/27 for all x  0 .\]
-\par
Given that each of the numbers a, b and c lies between 0 and 1, prove by contradiction\par
that at least one of the numbers bc(1 a), ca(1 b) and ab(1 c) is less than or equal\par
- - -\par
\[to 4/27.\]
\textbf{(ii)} Given that each of the numbers p and q lies between 0 and 1, prove that at least one of\par
\[the numbers p(1 q) and q(1 p) is less than or equal to 1/4.\]
- -\par
\end{flushleft}
\subsection*{Page 2}
\begin{flushleft}
\bigskip\textbf{Question 4.} A curve is given by\par
\[x^2 + y^2 + 2axy = 1,\]
where a is a constant satisfying 0 < a < 1. Show that the gradient of the curve at the point P\par
with coordinates (x, y) is\par
x + ay\par
,\par
- ax + y\par
\[provided ax + y = 0. Show that θ, the acute angle between OP and the normal to the curve\]
\par
at P , satisfies\par
\[tan θ = a y^2 x^2 .\]
| - |\par
dθ\par
\[Show further that, if = 0 at P , then:\]
dx\par
\textbf{(i)} a(x2 + y2) + 2xy = 0 ;\par
\textbf{(ii)} (1 + a)(x2 + y2 + 2xy) = 1 ;\par
a\par
\textbf{(iii)} tan θ = .\par
\[\sqrt\{1\} a^2\]
-\par
5 Evaluate the integrals\par
\[\pi/2 sin 2x \pi/2 sin x\]
dx and dx .\par
\[1 + sin2 x 1 + sin2 x\]
0 0\par
Show, using the binomial expansion, that (1 + \sqrt2 )5 < 99. Show also that \sqrt2 > 1.4. Deduce\par
that 2\sqrt2 > 1 + \sqrt2 . Use this result to determine which of the above integrals is greater.\par
\bigskip\textbf{Question 6.} A curve has the equation y = f(x), where\par
\[\pi 3x \pi\]
\[f(x) = cos 2x + + sin .\]
3 2 - 4\par
\textbf{(i)} Find the period of f(x).\par
\textbf{(ii)} Determine all values of x in the interval \pi  x  \pi for which f(x) = 0. Find a value\par
-\par
of x in this interval at which the curve touches the x-axis without crossing it.\par
\textbf{(iii)} Find the value or values of x in the interval 0  x  2\pi for which f(x) = 2 .\par
[Turn over\par
 \par
   \par
\end{flushleft}
\subsection*{Page 3}
\begin{flushleft}
\bigskip\textbf{Question 7.} (i) By writing y = u(1 + x2) 2 , where u is a function of x, find the solution of the equation\par
1 dy x\par
\[= xy +\]
\[y dx 1 + x^2\]
\[for which y = 1 when x = 0.\]
\textbf{(ii)} Find the solution of the equation\par
\[1 dy x^2\]
\[= x2y +\]
\[y dx 1 + x^3\]
\[for which y = 1 when x = 0.\]
\textbf{(iii)} Give, without proof, a conjecture for the solution of the equation\par
1 dy xn 1\par
\[= xn 1y + -\]
-\par
y dx 1 + xn\par
\[for which y = 1 when x = 0, where n is an integer greater than 1.\]
\bigskip\textbf{Question 8.} The points A and B have position vectors a and b, respectively, relative to the origin O. The\par
points A, B and O are not collinear. The point P lies on AB between A and B such that\par
\[AP : P B = (1 λ) : λ .\]
-\par
Write down the position vector of P in terms of a, b and λ. Given that OP bisects ∠AOB,\par
\[determine λ in terms of a and b, where a = a and b = b .\]
| | | |\par
\[The point Q also lies on AB between A and B, and is such that AP = BQ. Prove that\]
\[OQ2 OP 2 = (b a)2 .\]
- -\par
\end{flushleft}
\subsection*{Page 4}
\begin{flushleft}
\bigskip\textbf{Section B: Mechanics}\par
\bigskip\textbf{Question 9.} In this question, use g = 10 m s 2.\par
-\par
In cricket, a fast bowler projects a ball at 40 m s 1 from a point h m above the ground, which\par
-\par
is horizontal, and at an angle α above the horizontal. The trajectory is such that the ball will\par
strike the stumps at ground level a horizontal distance of 20 m from the point of projection.\par
\textbf{(i)} Determine, in terms of h, the two possible values of tan α.\par
Explain which of these two values is the more appropriate one, and deduce that the ball\par
hits the stumps after approximately half a second.\par
\textbf{(ii)} State the range of values of h for which the bowler projects the ball below the horizontal.\par
\textbf{(iii)} In the case h = 2.5, give an approximate value in degrees, correct to two significant\par
figures, for α. You need not justify the accuracy of your approximation.\par
[You may use the small-angle approximations cos θ 1 and sin θ θ.]\par
\approx \approx\par
\bigskip\textbf{Question 10.} The lengths of the sides of a rectangular billiards table ABCD are given by AB = DC = a\par
\[and AD = BC = 2b. There are small pockets at the midpoints M and N of the sides AD and\]
BC, respectively. The sides of the table may be taken as smooth vertical walls.\par
A small ball is projected along the table from the corner A. It strikes the side BC at X, then\par
the side DC at Y and then goes directly into the pocket at M . The angles BAX, CXY and\par
DYM are α, β and γ respectively. On each stage of its path, the ball moves with constant\par
speed in a straight line, the speeds being u, v and w respectively. The coefficient of restitution\par
\[between the ball and the sides is e, where e > 0.\]
\textbf{(i)} Show that tan α tan β = e and find γ in terms of α.\par
(1 + 2e)b\par
\textbf{(ii)} Show that tan α = and deduce that the shot is possible whatever the value\par
(1 + e)a\par
of e.\par
\textbf{(iii)} Find an expression in terms of e for the fraction of the kinetic energy of the ball that is\par
lost during the motion.\par
[Turn over\par
\end{flushleft}
\subsection*{Page 5}
\begin{flushleft}
\bigskip\textbf{Question 11.} A wedge of mass km has the shape (in cross-section) of a right-angled triangle. It stands on\par
a smooth horizontal surface with one face vertical. The inclined face makes an angle θ with\par
the horizontal surface. A particle P , of mass m, is placed on the inclined face and released\par
from rest. The horizontal face of the wedge is smooth, but the inclined face is rough and the\par
coefficient of friction between P and this face is µ.\par
\textbf{(i)} When P is released, it slides down the inclined plane at an acceleration a relative to the\par
wedge. Show that the acceleration of the wedge is\par
a cos θ\par
.\par
k + 1\par
To a stationary observer, P appears to descend along a straight line inclined at an\par
angle 45 to the horizontal. Show that\par
\textasciicircum{}\circ\par
k\par
\[tan θ = .\]
k + 1\par
\[In the case k = 3, find an expression for a in terms of g and µ.\]
\textbf{(ii)} What happens when P is released if tan θ  µ?\par
\end{flushleft}
\subsection*{Page 6}
\begin{flushleft}
\bigskip\textbf{Section C: Probability and Statistics}\par
\bigskip\textbf{Question 12.} In the High Court of Farnia, the outcome of each case is determined by three judges: the\par
ass, the beaver and the centaur. Each judge decides its verdict independently. Being simple\par
creatures, they make their decisions entirely at random. Past verdicts show that the ass gives\par
a guilty verdict with probability p, the beaver gives a guilty verdict with probability p/3 and\par
\[the centaur gives a guilty verdict with probability p^2.\]
Let X be the number of guilty verdicts given by the three judges in a case. Given that\par
\[E(X) = 4/3, find the value of p.\]
The probability that a defendant brought to trial is guilty is t. The King pronounces that the\par
defendant is guilty if at least two of the judges give a guilty verdict; otherwise, he pronounces\par
the defendant not guilty. Find the value of t such that the probability that the King pronounces\par
\[correctly is 1/2.\]
13 Bag P and bag Q each contain n counters, where n  2. The counters are identical in shape\par
and size, but coloured either black or white. First, k counters (0  k  n) are drawn at\par
random from bag P and placed in bag Q. Then, k counters are drawn at random from bag Q\par
and placed in bag P .\par
\textbf{(i)} If initially n 1 counters in bag P are white and one is black, and all n counters in\par
-\par
bag Q are white, find the probability in terms of n and k that the black counter ends\par
up in bag P .\par
Find the value or values of k for which this probability is maximised.\par
\textbf{(ii)} If initially n 1 counters in bag P are white and one is black, and n 1 counters in\par
- -\par
bag Q are white and one is black, find the probability in terms of n and k that the black\par
counters end up in the same bag.\par
Find the value or values of k for which this probability is maximised.\par
\end{flushleft}
\end{document}
